\section{Introduction}
\label{sec:introduction}

Whole-body vibration in the \SIrange{4}{8}{\hertz} range causes
gastrointestinal disturbances.  This is an established occupational-health
effect, codified in ISO~2631~\cite{iso2631}, reviewed by
Griffin~\cite{Griffin1990} and Mansfield~\cite{mansfield2005human},
and documented among truck drivers, helicopter pilots, and heavy-equipment
operators worldwide.  Airborne sound at the same frequencies does not produce
these effects, even at extreme intensities.  NASA-funded exposures at
140--154\,dB\,SPL caused chest-wall vibration and respiratory disruption, but
\emph{no gastrointestinal response}~\cite{Mohr1965}.  Two
excitation pathways, the same target organ, the same frequency band---yet
outcomes that differ by orders of magnitude.  This contrast motivates the
present study.

The ``brown note'' hypothesis---a hypothetical infrasonic tone alleged to
induce involuntary bowel evacuation---has also circulated in military
folklore, audio engineering, and popular media.  Scientifically,
the claim reflects a persistent confusion between what mechanical vibration
demonstrably does and what airborne sound is assumed to do.  A more plausible
airborne route, insofar as one exists, is
vestibular rather than visceral: recent measurements show that infrasound can
be transmitted through the ear strongly enough to engage acoustic and
vestibular pathways long before any whole-cavity abdominal resonance is
appreciably driven~\cite{Raufer2018}.

The empirical asymmetry between mechanical and airborne excitation is, in
fact, a straightforward consequence of long-wavelength acoustics.  At
the $n{=}2$ resonance ($f_2 \approx \SI{4}{\hertz}$) the acoustic wavelength
exceeds \SI{85}{\metre}---more than 250 times the abdominal diameter.  In
this deep Rayleigh regime, the $n$-th spatial harmonic of the incident field
scales as $(ka)^n$, where $ka \approx 0.01$.  Flexural deformation of the
abdominal wall ($n \geq 2$) therefore receives a driving pressure four orders
of magnitude below the free-field value.  The air--tissue impedance mismatch
compounds the deficit from the energy side: at the interface, approximately
\SI{99.9}{\percent} of incident acoustic power is reflected even before the
much smaller mode-specific resonant absorption is considered, so the energy
available to drive flexural modes is negligible.  Mechanical
vibration, transmitted through the skeleton and abdominal wall, suffers
neither penalty.  The ratio between these two pathways---the \emph{coupling
disparity}---is the central quantity we derive and the principal contribution
of this paper.

Scientific discussion of infrasound predates later popular-media portrayals,
including the \emph{South Park} episode ``World Wide Recorder Concert'' (also
known as ``The Brown Noise'')~\cite{SouthPark2000BrownNoise}.  Gavreau's
experiments at CNRS Levallois-Perret
in the late 1950s, where a faulty ventilator at approximately
\SI{7}{\hertz} caused nausea among laboratory
staff~\cite{Gavreau1966,Leventhall2007},
drew early attention to infrasound.  Von~Gierke's systematic review for
NASA~\cite{vonGierke1974,vonGierke2002} established exposure criteria that
informed later ISO standards.  Tandy and Lawrence~\cite{Tandy1998}
demonstrated that a \SI{19}{\hertz} standing wave could induce visual
disturbances and anxiety---real perceptual effects, but far from the claimed
bowel response.  Critical reviews by
Leventhall~\cite{Leventhall2003,Leventhall2007} and assessments of acoustic
weapons by Altmann~\cite{Altmann1999} and Jauchem and
Cook~\cite{Jauchem2007} concluded that infrasound is impractical as a weapon
and that claims of extreme physiological effects are exaggerated or
fabricated.  Recent wind-turbine reviews likewise find that the dominant
health concerns are annoyance, sleep disturbance, and other indirect effects,
not visceral resonance~\cite{vanKamp2021}.  Randomised exposure studies with
neuroimaging have also reported subtle but measurable neural correlates of
airborne infrasound exposure, but again no evidence of gastrointestinal or
abdominal-resonance effects~\cite{Ascone2021}.  Yet none of this work quantified \emph{why} airborne and
mechanical pathways differ so dramatically---the conclusion was empirical,
not mechanistic.

We provide the mechanistic account.  To our knowledge, no prior analytical study
has combined first-principles cavity--shell modal analysis with an
energy-consistent comparison of airborne and mechanical coupling pathways:
\begin{enumerate}
  \item derived abdominal cavity eigenfrequencies from first-principles
    shell theory with proper fluid--structure coupling;
  \item quantified the airborne acoustic coupling coefficient for flexural
    modes in the long-wavelength limit; or
  \item compared airborne and mechanical excitation pathways within a common
    energy-budget framework to obtain a single, quantitative coupling ratio.
\end{enumerate}

We address all three gaps.  We model the abdomen as a fluid-filled
viscoelastic oblate spheroidal shell, derive natural frequencies for both
breathing ($n=0$) and flexural ($n \geq 2$) mode families, and compare
airborne acoustic coupling---subject to a severe $(ka)^n$ penalty at
infrasonic wavelengths---against direct mechanical (whole-body) vibration
coupling, which suffers no such penalty.  More broadly, this provides a
general template for comparing excitation pathways in acoustically compact,
fluid-loaded compliant structures.  The same framework identifies when
geometry, impedance mismatch, and modal participation make one forcing route
dominant over another even when both act on the same target mode.

The headline result is an upper-bound coupling disparity of approximately
$3.3 \times 10^4$ between mechanical and airborne pathways for
\SI{120}{\decibel} SPL airborne excitation and
$a_\mathrm{rms} = \SI{0.1}{\metre\per\second\squared}$ RMS whole-body
vibration under the SDOF idealisation ($\Gamma_2 = 1$).  Including the
canonical participation factor $\Gamma_2 \approx 0.48$ reduces this to
$\sim 1.6 \times 10^4$.  This qualitative ordering persists across
physiologically realistic parameter ranges, as we confirm through Monte
Carlo uncertainty quantification, and resolves the empirical asymmetry
cleanly.  It
explains why whole-body vibration at \SIrange{4}{8}{\hertz} produces
well-documented gastrointestinal effects, why airborne infrasound at
comparable frequencies does not, and why claims of a brown note are not
supported by the physics of airborne acoustic coupling.

The remainder of this paper is organised as follows.
Section~\ref{sec:formulation} develops the mathematical formulation.
Section~\ref{sec:results} presents a parametric study and compares the model
against ISO~2631 data.  Section~\ref{sec:coupling} quantifies the coupling
disparity.  Section~\ref{sec:discussion} discusses implications, alternative
pathways, and limitations, and Section~\ref{sec:conclusion} draws conclusions.
